\chapter{工程分析}

\section{建设项目概况}

\begin{table}[H]
\centering
\begin{tabular}{|c|c|}
\hline
项目 & 内容 \\
\hline
项目名称 & \projectnamefull \\
\hline
建设单位 & \buildunit \\
\hline
法定代表人 & \buildunitlegalrep \\
\hline
联系人 & \buildunitcontact \\
\hline
联系电话 & \buildunitphone \\
\hline
建设地址 & \buildaddress \\
\hline
项目性质 & \projectnature \\
\hline
建设规模 & \constructionscale \\
\hline
设备配置 & \equipmentconfig \\
\hline
总投资 & \totalinvestment \\
\hline
放射性危害程度 & \radiationhazarddegree \\
\hline
诊疗风险危害分类 & \textbf{\riskclassification} \\
\hline
设计单位 & \designunit \\
\hline
\end{tabular}
\end{table}


\section{人员结构}


\begin{table}[H]
\centering
\caption{现有在岗放射工作人员}
\begin{tabular}{|c|c|c|c|c|c|c|}
\hline
序号 & 姓名 & 学历 & 专业 & 职称 & 资质证书 & 岗位 \\
\hline
1 & 陈明伟 & 本科 & 肿瘤内科 & 中级 & LA 物理师证 & 医师 \\
\hline
2 & 刘鲁根 & 本科 & 临床医学 / 工程 & 中级 & LA 物理师证 & 物理师 \\
\hline
3 & 张红宇 & 大专 & 放射技术 & 中级 & LA 技师证 & 技师 \\
\hline
4 & 幸焕中 & 本科 & 放射医学技术 & 中级 & LA 技师证 & 技师 \\
\hline
5 & 欧阳公怒 & 本科 & 临床医学工程 & 中级 & 无 & 设备辅助维修 \\
\hline
\end{tabular}
\end{table}



\begin{table}[H]
\centering
\caption{项目人员配置合规性}
\begin{tabular}{|c|c|c|c|c|}
\hline
岗位配置 & 现有人数 & 新增 & 法规要求 & 结论 \\
\hline
放疗医师 & 1 & 无 & 中级以上放疗医师 & 符合 \\
\hline
医学物理师 & 1 & 无 & 本科 / 中级以上物理人员 & 符合 \\
\hline
放疗技师 & 2 & 无 & 持证放疗技师 & 符合 \\
\hline
设备维修 & 1 & 无 & 设备维修人员 & 符合 \\
\hline
\end{tabular}
\end{table}



\begin{table}[H]
\centering
\caption{岗位职责}
\begin{tabular}{|c|c|c|}
\hline
岗位 & 定员 & 工作职责 \\
\hline
放疗医师 & 1 & 靶区勾画、治疗方案制定、质控管理 \\
\hline
物理师 & 1 & 放疗计划设计、剂量核算 \\
\hline
放疗技师 & 2 & 设备操作、患者摆位、日常质控 \\
\hline
维修 & 1 & 设备日常检修维护 \\
\hline
\end{tabular}
\end{table}


\section{周围环境}

项目位于医院东北侧肿瘤大楼一层；
东侧：常青路；南侧：直冲小区 + 城市绿地；西侧：黄塘巷；北侧：武功山中大道；
机房南侧为预留直线加速器机房，西侧为控制室、准备间，东侧为狭窄过道，上方屋顶平台、下方素土地层。

\section{环境本底辐射}

室内本底：83\textbackslash\{\}\textbackslash\{\}92nGy/h；室外：33\textbackslash\{\}\textbackslash\{\}65nGy/h；
数据符合萍乡天然本底范围（室内 64.7\textbackslash\{\}\textbackslash\{\}149.3nGy/h、室外 14.4\textbackslash\{\}\textbackslash\{\}168.8nGy/h）。

\section{项目工艺与原理}

\subsection{后装机原理}

近距离腔内放疗，$^{192}\text{Ir}$（半衰期 74d，γ 平均能量 0.37MeV）；
非治疗时源密闭贮源器：5cm 处≤50μSv/h，1m 处≤5μSv/h；
治疗流程：患者置入施源器→CT 模拟定位→TPS 制定计划→转运至机房→工作人员外接管路→人员撤离关门→远程送源照射→治疗结束源自动回储源器。

\subsection{完整诊疗流程}

患者就诊→适应症评估（正当性）→预约→准备间置施源器→CT 机房定位扫描→物理师做治疗计划→转入后装机房→管路接驳→工作人员离场→遥控治疗→结束撤管、患者离院。

\section{预期运行工作量}

\begin{table}[H]
\centering
\begin{tabular}{|c|c|c|c|c|c|c|}
\hline
设备 & 日接诊 & 单次照射时长 & 周工作日 & 年工作周 & 周总工时 & 年总工时 \\
\hline
后装治疗机 & 6 人次 & 10min & 5d & 50 周 & 5.0h & 250h \\
\hline
\end{tabular}
\end{table}


\begin{quote}
备注：每周质控 10min，全年合计工作 254h。
\end{quote}

\section{利旧内容}

\begin{enumerate}
\item \textbf{人员利旧}：全部沿用现有持证放射工作人员；
\item \textbf{设备利旧}：新购主机，部分质控、监测仪器利旧；
\item \textbf{建筑利旧}：原有机房墙体、混凝土屏蔽、防护铅门全部利旧；
\item \textbf{管理利旧}：原有放射管理制度、应急预案、人员健康 / 剂量 / 档案管理体系沿用并完善。
\end{enumerate}

\section{小结}

本项目为改建、\radiationhazarddegree 放射项目；单台设备配$^{192}\text{Ir}$Ⅲ 类密封源；人员配置满足法规要求，原有建筑主体保留。
